\chapter{Συμπεράσματα}
\label{ch:conclusions}

Στο κεφάλαιο αυτό συνοψίζονται τα κύρια ευρήματα της εργασίας, συζητούνται
οι περιορισμοί και προτείνονται κατευθύνσεις για μελλοντική έρευνα.

% ─────────────────────────────────────────────────────────────────────────────
\section{Σύνοψη Κύριων Ευρημάτων}
\label{sec:conc_findings}
% ─────────────────────────────────────────────────────────────────────────────

Η παρούσα εργασία διερεύνησε συστηματικά τέσσερις στρατηγικές βραχυπρόθεσμης
πρόβλεψης (Teacher-Forcing, Recursive/Open-Loop, MIMO και Walk-Forward Monthly
Retraining) για την τιμή DAM και το ηλεκτρικό φορτίο στην ελληνική αγορά,
εφαρμόζοντας πέντε αλγόριθμους ML σε τρεις χρονικές περιόδους αξιολόγησης.

\subsection*{ΕΕ1: Ποια στρατηγική είναι καλύτερη;}

\textbf{Τιμή DAM:} Η Walk-Forward Monthly Retraining (TF) αναδεικνύεται
ξεκάθαρα η βέλτιστη στρατηγική για μακροχρόνιες περιόδους:
\begin{itemize}
    \item Q1 2026: MR Ensemble-TF-Best3 = 10.491~€/MWh vs
    static CL Ensemble = 14.478~€/MWh (\textbf{−27,7\%})
    \item Η υπεροχή επιβεβαιώνεται στατιστικά σημαντικά για όλα τα ζεύγη
    (DM test, $p < 0{,}05$).
\end{itemize}

\textbf{Φορτίο:} Η MR-TF και το static CL Ensemble είναι στατιστικά
αδιάκριτα (DM test, $p = 0{,}12$). Πρακτικά ωστόσο, η MR-TF δίνει
MAE 67,7~MW vs 96,8~MW (−30\%) — σημαντικό πρακτικά, αν και όχι
πάντα στατιστικά σημαντικό.

\subsection*{ΕΕ2: Πώς αλλάζει η κατάταξη με τον ορίζοντα;}

\begin{itemize}
    \item \textbf{Βραχυχρόνια (1 εβδομάδα):} CL/TF = 9--11 €/MWh,
    OL = 15--16 €/MWh, MIMO = 25--26 €/MWh. Το CL κυριαρχεί σαφώς.
    \item \textbf{Μακροχρόνια (3 μήνες):} CL MAE αυξάνει σε 14,5 €/MWh
    λόγω distributional shift. Η MR «κλειδώνει» MAE ≈ 10--11 €/MWh/μήνα
    με μηνιαία επαναδιαμόρφωση. Για Q4 2025 (υψηλή μεταβλητότητα Οκτ./Νοε.),
    ακόμα και το Naive-1 (19.5 €/MWh) νικά τα weekly OL μοντέλα.
    \item \textbf{OL daily (H=24) vs weekly (H=168):} H=24 chunks μειώνει
    OL MAE κατά 16,7\% για Δεκέμβριο 2025 (14,5 vs 17,4 €/MWh), εξαλείφοντας
    error accumulation 7 ημερών.
\end{itemize}

\subsection*{ΕΕ3: Στατιστική σημαντικότητα διαφορών;}

Ο DM test επιβεβαιώνει:
\begin{itemize}
    \item Walk-Forward TF > όλα (τιμή, $p<0{,}05$)
    \item CL ≈ Walk-Forward TF για φορτίο (μη σημαντική διαφορά)
    \item Εντός-στρατηγικής LGBM ≈ XGB ≈ RF (μη σημαντικές διαφορές)
    \item MIMO < CL/TF για αμφότερες τις εργασίες ($p<0{,}05$)
\end{itemize}

\subsection*{ΕΕ4: Ποια χαρακτηριστικά αξιοποιούν τα μοντέλα;}

\begin{itemize}
    \item \textbf{TF:} Κυριαρχεί $y_{t-1}$ (88\% για LGBM) — αυτό-αναδρομική δυναμική.
    \item \textbf{Recursive:} Μικρότερη εξάρτηση από $y_{t-1}$ (41\%),
    αυξημένη βαρύτητα $y_{t-24}$, $y_{t-168}$.
    \item \textbf{MIMO:} Ανάδυση εξωγενών (gas\_price, αιολική παραγωγή);
    αναγκαία στροφή λόγω μη πρόσβασης σε πραγματικές ενδιάμεσες τιμές.
    \item \textbf{Dense lags (MIMO):} Αύξηση συλλογικής βαρύτητας lags 4--23,
    οδηγώντας σε βελτίωση MLP −22\%, XGB −7,9\%, LGBM −6,5\%.
\end{itemize}

\subsection*{ΕΕ5: Αντισταθμίζει η MR το distributional shift;}

\textbf{Ναι, και εντυπωσιακά.} Το σημαντικότερο εύρημα της εργασίας:

\begin{table}[htbp]
\centering
\caption{Σύγκριση στρατηγικών Q1 2026 — κύρια αποτελέσματα.}
\label{tab:conc_summary}
\begin{tabular}{lrr}
\toprule
\textbf{Στρατηγική} & \textbf{Τιμή (€/MWh)} & \textbf{Φορτίο (MW)} \\
\midrule
Walk-Forward MR (TF Ensemble)     & \textbf{10,491} & \textbf{67,729} \\
Static CL (TF Ensemble)           & 14,478           & 96,751 \\
Static OL — best single (H=24)    & 20,094           & 102,961 \\
MIMO — best ensemble              & 20,896           & 217,088 \\
\midrule
Naive-1 (baseline)                & 15,433           & 267,151 \\
\bottomrule
\end{tabular}
\end{table}

Η μηνιαία επανεκπαίδευση αντισταθμίζει πλήρως τη μεταβολή κατανομής
(Φεβρουάριος 2026: μέση τιμή 78~€/MWh έναντι 110~€/MWh Ιανουαρίου): το
TF-LGBM-MR επιτυγχάνει MAE 11,7~€/MWh Φεβρουαρίου, έναντι 15,9~€/MWh
της static CL.

% ─────────────────────────────────────────────────────────────────────────────
\section{Ειδικά Ευρήματα}
\label{sec:conc_special}
% ─────────────────────────────────────────────────────────────────────────────

\paragraph{Dense lags paradox.}
Τα dense lags (4--23) βελτιώνουν MIMO (χωρίς recursion $\to$ καθαρές ιστορικές τιμές)
αλλά \textit{βλάπτουν} το Recursive OL (+2,8\% LGBM). Αιτία: στο OL, οι
dense lags είναι αυτο-αναφορικές προβλέψεις με συσσωρευμένο σφάλμα. Το Scheduled
Sampling αντισταθμίζει μερικά αυτό το πρόβλημα (XGB-SS-Dense βελτιώνεται +1,4--1,5\%).

\paragraph{SVR κατάρρευση.}
Το SVR αποτυγχάνει εντυπωσιακά για Q1 2026 (τιμή: 40,8~€/MWh, φορτίο: 600~MW).
Αιτία: η βαθμονόμηση του πυρήνα RBF (trained on 2017--Nov2025) δεν γενικεύεται
στη νέα κατανομή Ιανουαρίου--Φεβρουαρίου 2026 (volatile, low-price regime).
Το SVR είναι ευαίσθητο στη μεταβολή κατανομής λόγω μη-παραμετρικής φύσης.

\paragraph{Ensemble ανθεκτικότητα.}
Το Ensemble-Best3 (1/MAE weights) αποδεικνύεται σταθερά ανθεκτικό: για φορτίο
MR, Ensemble-TF-Best3 = 67,7~MW vs TF-LGBM = 69,4~MW, TF-XGB = 72,9~MW.
Ακόμα και όταν ένα μοντέλο (RF ή MLP) αποδίδει άσχημα σε συγκεκριμένο μήνα,
το ensemble περιορίζει τη ζημιά λόγω χαμηλής βαρύτητας.

\paragraph{Scheduled Sampling για φορτίο.}
Για φορτίο, το Scheduled Sampling (Rec-XGB-SS-DO) επιτυγχάνει 70,6~MW — σχεδόν
ισοδύναμο με TF-Ensemble (67,7~MW). Αυτό αποδεικνύει ότι για σήματα με ομαλή
δυναμική (φορτίο), η Recursive στρατηγική + SS μπορεί να ανταγωνιστεί την
oracle TF.

% ─────────────────────────────────────────────────────────────────────────────
\section{Περιορισμοί}
\label{sec:conc_limitations}
% ─────────────────────────────────────────────────────────────────────────────

\begin{enumerate}
    \item \textbf{Συνθετικές τιμές φυσικού αερίου:} Για την περίοδο Νοέμβριος
    2025 -- Φεβρουάριος 2026, οι τιμές TTF gas υπολογίστηκαν με γραμμική
    παρεμβολή από σημεία αγκύρωσης. Ο πραγματικός TTF Ιανουαρίου 2026
    κυμάνθηκε σε 42--50~€/MWh, ενώ οι συνθετικές τιμές εκτιμούσαν ~39~€/MWh.
    Αυτό εισάγει θόρυβο χαρακτηριστικών, επηρεάζοντας κυρίως τα μοντέλα
    τιμής Ιανουαρίου--Φεβρουαρίου 2026.

    \item \textbf{Subsampling SVR:} Το SVR εκπαιδεύτηκε σε μόνο 5.000 δείγματα
    (~7 μήνες δεδομένων, από συνολικά ~80.000) λόγω υπολογιστικού κόστους
    $O(n^2)$--$O(n^3)$. Τα αποτελέσματα SVR δεν είναι αντιπροσωπευτικά.

    \item \textbf{Απουσία θερμοκρασιακών δεδομένων:} Η θερμοκρασία αποτελεί
    κρίσιμο χαρακτηριστικό για ELF. Η παρούσα εργασία αντιμετωπίζει αυτό μέσω
    rolling averages και ημερολογιακών χαρακτηριστικών ως proxy, αλλά η
    άμεση ενσωμάτωση θερμοκρασιακής πρόγνωσης αναμένεται να βελτιώσει
    την ακρίβεια φορτίου.

    \item \textbf{Μέγεθος δοκιμαστικής περιόδου:} Η αξιολόγηση σε 3 μήνες
    Q1 2026 παρέχει εύλογο μέγεθος δείγματος ($T = 2160$) αλλά αντιστοιχεί
    σε ένα μόνο τρίμηνο. Η γενίκευση αποτελεσμάτων σε ολόκληρο έτος ή σε
    διαφορετικές εποχές παραμένει υποθετική.

    \item \textbf{Αθροιστική συγχώνευση γεννήτριας δεδομένων:} Η μετατροπή
    15-λεπτων αναγνώσεων παραγωγής σε ωριαία χρησιμοποιεί SUM αντί MEAN,
    αποτελώντας γνωστό τεχνικό παρατηρούμενο. Δεδομένου ότι η κλίμακα
    παραμένει σταθερή μεταξύ εκπαίδευσης και δοκιμής, η επίδραση στην
    ακρίβεια πρόβλεψης είναι μηδενική.
\end{enumerate}

% ─────────────────────────────────────────────────────────────────────────────
\section{Μελλοντική Έρευνα}
\label{sec:conc_future}
% ─────────────────────────────────────────────────────────────────────────────

\begin{enumerate}
    \item \textbf{Θερμοκρασιακά δεδομένα:} Ενσωμάτωση ωριαίων μετεωρολογικών
    προβλέψεων από ECMWF/GFS (θερμοκρασία, ηλιοφάνεια, ταχύτητα ανέμου) ως
    χαρακτηριστικά — αναμένεται σημαντική βελτίωση ιδιαίτερα για ELF.

    \item \textbf{Βαθιά Μάθηση (LSTM/Transformer):} Δοκιμή αρχιτεκτονικών
    ακολουθιών (seq2seq LSTM, Temporal Fusion Transformer~\cite{lim2021temporal})
    που διαχειρίζονται εγγενώς αναδρομικές δομές, πιθανώς υπερεκδίδοντας τα
    τρέχοντα MIMO/Recursive ML μοντέλα.

    \item \textbf{Online Learning:} Αντί μηνιαίας επανεκπαίδευσης, διερεύνηση
    αλγόριθμων online/incremental learning (π.χ. gradient boosting με
    νέα δεδομένα κάθε ημέρα) που προσαρμόζονται συνεχώς χωρίς πλήρη
    επανεκπαίδευση.

    \item \textbf{Πιθανοτική Πρόβλεψη:} Επέκταση σε probabilistic forecasting
    (prediction intervals, quantile regression) για παροχή εκτιμήσεων
    αβεβαιότητας — κρίσιμης αξίας για εμπόρους ενέργειας.

    \item \textbf{Two-Stage Load$\to$Price:} Αν και η τρέχουσα εργασία έδειξε
    αμελητέα επίδραση (load\_fc δεν είναι κρίσιμο χαρακτηριστικό τιμής),
    η διερεύνηση σε περιόδους υψηλής αβεβαιότητας φορτίου μπορεί να
    αναδείξει διαφορετική εικόνα.

    \item \textbf{Γενίκευση σε άλλες αγορές:} Επέκταση της μεθοδολογίας σε
    άλλες ευρωπαϊκές αγορές (Βουλγαρία, Ρουμανία, Σερβία) που εμφανίζουν
    παρόμοια χαρακτηριστικά (υψηλή διείσδυση ΑΠΕ, διασυνδέσεις, ευμετάβλητες
    τιμές).
\end{enumerate}

% ─────────────────────────────────────────────────────────────────────────────
\section{Κλείσιμο}
\label{sec:conc_closing}
% ─────────────────────────────────────────────────────────────────────────────

Η εργασία αυτή καταδεικνύει ότι η \textit{επιλογή στρατηγικής πρόβλεψης} έχει
εξίσου ή και μεγαλύτερη σημασία από την επιλογή αλγόριθμου ML: η Walk-Forward
Monthly Retraining βελτιώνει κατά 28--30\% έναντι static training, ανεξαρτήτως
αλγόριθμου. Η ορθή μεθοδολογική αντιμετώπιση — αποφυγή data leakage, πλήρης
out-of-sample αξιολόγηση, στατιστικοί έλεγχοι — αποτελεί τη βάση αξιόπιστης
σύγκρισης στη βιβλιογραφία EPF/ELF.

Τα αποτελέσματα έχουν άμεση πρακτική σημασία για τους συμμετέχοντες στην
ελληνική αγορά ηλεκτρισμού: αναλυτές ρίσκου, εμπόρους ενέργειας και Διαχειριστή
Συστήματος Μεταφοράς που αντιμετωπίζουν τις ίδιες προκλήσεις μεταβλητότητας
και distributional shift.
